% Part: first-order-logic
% Chapter: introduction
% Section: substitution

\documentclass[../../../include/open-logic-section]{subfiles}

\begin{document}

\olfileid{fol}{int}{sub}

\olsection{Substitutie}

Een voorbeeld laat zien hoe alles samenhangt. Het laat ook zien waarom
we syntaxis en semantiek zo precies opbouwen: die vormen de basis voor
het verdere onderzoek. Met onze !!{derivation}ssystemen moeten we de
formule $\Atom{\Obj P}{\Obj a}$ uit
$\lforall[\Obj v_0][\Atom{\Obj P}]{\Obj v_0}$ kunnen !!{derive}. Misschien willen we daar zelfs een
afleidingsregel van maken. Dan moet de regel wel algemeen gelden. Hij
mag niet alleen over $\Obj P$, $\Obj a$ en $\Obj v_0$ gaan, maar moet
werken voor iedere !!{formula}~$!A$, iedere term~$t$ en iedere
!!{variable}~$x$. (Een !!{constant} is een term. We bekijken later ook
ingewikkeldere termen die zijn opgebouwd uit !!{constant}s en
!!{function}s.) We willen dus kunnen zeggen: ``Kun je
$\lforall[x][!A(x)]$ !!{derive}, dan mag je~$!A(t)$ afleiden. Je haalt
$\lforall[x]$ weg en vervangt $x$ door~$t$.'' Maar wat houdt
``$x$ door~$t$ vervangen'' precies in? Hoe hangen $!A(x)$ en~$!A(t)$
samen? En werkt dit altijd?

We maken dit precies met de bewerking \emph{substitutie}. Die bewerking
is lastiger dan zij lijkt. We mogen niet zomaar alle voorkomens van~$x$
in~$!A$ door~$t$ vervangen. Ook mag niet iedere~$t$ voor iedere~$x$
worden ingevuld. We lossen dit opnieuw op met inductieve definities.
Daarna is de regel ``leid $!A(t)$ af uit $\lforall[x][!A(x)]$'' wel
precies bepaald. We kunnen dan ook bewijzen dat het een goede regel is:
$\lforall[x][!A(x)]$ heeft~$!A(t)$ als semantisch gevolg.
Om te bewijzen dat $\lforall[x][!A(x)]$ semantisch~$!A(t)$ tot gevolg
heeft, is nog een feit nodig waarvan het nut niet meteen duidelijk
hoeft te zijn. Of $\Sat{M}{!A(t)}$ geldt, hangt alleen af van
de waarde van de term~$t$. Laat $m$ het !!{element} van~$\Domain{M}$
zijn dat $t$ aanwijst. Dan geldt $\Sat{M}{!A(t)}[s]$ precies als
$\Sat{M}{!A(x)}[\Subst{s}{m}{x}]$. Dit blijft gelden wanneer $t$
!!{variable}s bevat. We moeten dan wel heel precies zeggen wat de
stelling beweert.

\end{document}
